% ch4_a 引力 (书页 151-162 / p166-p177)
% 第4章 Gravitation
\chapter{引力}

\section{弯曲时空中的物理}

数学上的会费已经缴清，我们现在可以着手考察广义相对论所描述的引力物理了。这个课题自然地分成两部分：引力场如何影响物质的行为，以及物质如何决定引力场。在 Newton 引力中，这两部分内容分别是：物体在引力势 $\Phi$ 中加速度的表达式
\begin{equation*}
\mathbf{a} = -\nabla\Phi,
\tag{4.1}
\end{equation*}
以及用物质密度 $\rho$ 和 Newton 引力常数 $G$ 表出引力势的 Poisson 微分方程：
\begin{equation*}
\nabla^2\Phi = 4\pi G\rho.
\tag{4.2}
\end{equation*}

在广义相对论中，与之相应的表述将描述时空曲率如何作用于物质、使自身表现为引力，以及能量和动量如何影响时空、使之产生曲率。无论从哪一半入手，直奔主题都是说得通的：径直给出支配弯曲时空中物理的定律，再推导其后果。不过，我们想讲得更有动机引导一些：从基本物理原理出发，试图论证这些原理自然地导致一个几乎是唯一的物理理论。

在第 2 章中，我们通过引入 Einstein 等效原理（Einstein Equivalence Principle，简称 EEP）来引出对流形的讨论：“在足够小的时空区域内，物理定律退化为狭义相对论的定律；不可能借助局部实验探测到引力场的存在。”EEP 源于这样一个观念：引力是\emph{普适的}（universal）；它以同样的方式影响所有粒子（实际上包括一切形式的能量-动量）。正是这种普适性促使 Einstein 提出：我们所体验到的引力，是时空曲率的一种显现。这个想法很简单：像引力这般普适的东西，最容易描述为物质场传播所处背景的一种基本属性，而不是一种常规的力。与此同时，“时空是弯曲流形”这一认定还得到一个相似性的支持：一方面引力在局部区域中不可探测，另一方面我们却能在流形上找到局部惯性坐标系（在一点 $p$ 处 $g_{\hat\mu\hat\nu}=\eta_{\hat\mu\hat\nu}$、$\partial_{\hat\rho}g_{\hat\mu\hat\nu}=0$）。

最妙的是，这些抽象的哲学思辨可以直接转化为一条把物理定律推广到弯曲时空情形的简单处方，这就是\textbf{最小耦合原理}（minimal-coupling principle）。用最直白的方式来表述，这条处方可以陈述如下：

\begin{enumerate}
\item 取一条在平直时空惯性系坐标中成立的物理定律。
\item 把它写成坐标不变（张量）的形式。
\item 断言所得定律在弯曲时空中依然成立。
\end{enumerate}

把如此简单的想法铺陈成一个三步骤流程，也许显得有点小题大做。我们只希望说明白：这里头并没有什么太复杂的东西。在实际操作上，这条处方通常等价于：取一条公认的平直空间中的定律，把 Minkowski 度规 $\eta_{\mu\nu}$ 换成更一般的度规 $g_{\mu\nu}$，并把偏导数 $\partial_\mu$ 换成协变导数 $\nabla_\mu$。因此，在那些用逗号和分号分别记偏导数与协变导数的人中间，这条处方有时被称为“逗号变分号法则”（Comma-Goes-to-Semicolon Rule）。

举一个直截了当的例子：考虑自由下落（无加速度）粒子的运动。在平直空间中，这类粒子沿直线运动；用方程来表达，就是参数化路径 $x^\mu(\lambda)$ 的二阶导数为零：
\begin{equation*}
\frac{d^2x^\mu}{d\lambda^2} = 0.
\tag{4.3}
\end{equation*}
在一般坐标下，这并不是一个张量方程；虽然 $dx^\mu/d\lambda$ 是一个良定义矢量的分量，二阶导数分量 $d^2x^\mu/d\lambda^2$ 却不是。你可能真的会以为这是一条看上去像张量的方程；然而你可以很容易地验证，它在极坐标下甚至都不成立——除非你指望自由粒子做圆周运动。我们可以用链式法则写出
\begin{equation*}
\frac{d^2x^\mu}{d\lambda^2} = \frac{dx^\nu}{d\lambda}\,\partial_\nu\frac{dx^\mu}{d\lambda}.
\tag{4.4}
\end{equation*}
现在很清楚该如何把它推广到弯曲空间——只要把偏导数换成协变导数，
\begin{equation*}
\frac{dx^\nu}{d\lambda}\,\partial_\nu\frac{dx^\mu}{d\lambda}
\;\to\;
\frac{dx^\nu}{d\lambda}\,\nabla_\nu\frac{dx^\mu}{d\lambda}
= \frac{d^2x^\mu}{d\lambda^2} + \Gamma^\mu_{\rho\sigma}\frac{dx^\rho}{d\lambda}\frac{dx^\sigma}{d\lambda}.
\tag{4.5}
\end{equation*}
于是我们认识到，Newton 关系式 (4.3) 恰当的广义相对论版本就是测地线方程（geodesic equation），
\begin{equation*}
\frac{d^2x^\mu}{d\lambda^2} + \Gamma^\mu_{\rho\sigma}\frac{dx^\rho}{d\lambda}\frac{dx^\sigma}{d\lambda} = 0.
\tag{4.6}
\end{equation*}
因此在广义相对论中，自由粒子沿测地线运动；这一点我们此前已经提过，但现在你对它为什么成立应该有了稍好一点的理解。

作为一个还要更加直截了当的例子——而且我们已经提到过它——有平直时空中的能量-动量守恒定律：
\begin{equation*}
\partial_\mu T^{\mu\nu} = 0.
\tag{4.7}
\end{equation*}
代入我们的处方，就得到向弯曲时空的恰当推广：
\begin{equation*}
\nabla_\mu T^{\mu\nu} = 0.
\tag{4.8}
\end{equation*}
事情就是这么简单——以至于我们在第 3 章里使用这个方程时感到相当心安理得，并未给出任何细致的论证。当然，这种简单性丝毫不会减损向弯曲时空推广所带来的深刻后果，膨胀宇宙的例子便是明证。

把一个方程从平直时空推广到弯曲时空是一回事；而论证所得的结果描述的就是引力，则完全是另一回事。为此，我们可以证明 Newton 引力的通常结果如何纳入这幅图像。我们用三个要求来定义 Newton 极限（Newtonian limit）：粒子运动缓慢（相对于光速）、引力场很弱（从而可以视为对平直空间的微扰），而且场还是静态的（不随时间改变）。让我们取固有时 $\tau$ 作为仿射参量，看看这些假设会把测地线方程变成什么。“缓慢运动”意味着
\begin{equation*}
\frac{dx^i}{d\tau} \ll \frac{dt}{d\tau},
\tag{4.9}
\end{equation*}
于是测地线方程变为
\begin{equation*}
\frac{d^2x^\mu}{d\tau^2} + \Gamma^\mu_{00}\left(\frac{dt}{d\tau}\right)^2 = 0.
\tag{4.10}
\end{equation*}
由于场是静态的（$\partial_0 g_{\mu\nu}=0$），相关的 Christoffel 符号 $\Gamma^\mu_{00}$ 得以简化：
\begin{align*}
\Gamma^\mu_{00} &= \frac{1}{2}g^{\mu\lambda}\left(\partial_0 g_{\lambda 0} + \partial_0 g_{0\lambda} - \partial_\lambda g_{00}\right) \\
&= -\frac{1}{2}g^{\mu\lambda}\partial_\lambda g_{00}.
\tag{4.11}
\end{align*}
最后，引力场的微弱使我们能把度规分解成 Minkowski 形式加上一个小微扰：
\begin{equation*}
g_{\mu\nu} = \eta_{\mu\nu} + h_{\mu\nu}, \qquad |h_{\mu\nu}| \ll 1.
\tag{4.12}
\end{equation*}

我们采用的是惯性系坐标，因此 $\eta_{\mu\nu}$ 是度规的规范形式（canonical form）。对度规微扰 $h_{\mu\nu}$ 的“小量条件”在任意坐标系下其实并无意义。由逆度规的定义 $g^{\mu\nu}g_{\nu\sigma}=\delta^\mu_\sigma$，我们发现准确到 $h$ 的一阶，
\begin{equation*}
g^{\mu\nu} = \eta^{\mu\nu} - h^{\mu\nu},
\tag{4.13}
\end{equation*}
其中 $h^{\mu\nu} = \eta^{\mu\rho}\eta^{\nu\sigma}h_{\rho\sigma}$。事实上，对 $h$ 的任一确定阶数的对象，我们都可以用 Minkowski 度规来升降指标，因为修正项只会在更高阶才有贡献。如果你愿意，可以把 $h_{\mu\nu}$ 想成一个在 Minkowski 空间中传播、并与其他场相互作用的对称 $(0,2)$ 张量场。

把所有这些合在一起，准确到 $h_{\mu\nu}$ 的一阶，我们得到
\begin{equation*}
\Gamma^\mu_{00} = -\frac{1}{2}\eta^{\mu\lambda}\partial_\lambda h_{00}.
\tag{4.14}
\end{equation*}
因此测地线方程 (4.10) 成为
\begin{equation*}
\frac{d^2x^\mu}{d\tau^2} = \frac{1}{2}\eta^{\mu\lambda}\partial_\lambda h_{00}\left(\frac{dt}{d\tau}\right)^2.
\tag{4.15}
\end{equation*}
利用 $\partial_0 h_{00}=0$，它的 $\mu=0$ 分量就是
\begin{equation*}
\frac{d^2t}{d\tau^2} = 0.
\tag{4.16}
\end{equation*}
也就是说，$dt/d\tau$ 是常数。为了考察 (4.15) 的类空分量，回忆 $\eta^{\mu\nu}$ 的类空分量正是 $3\times 3$ 单位矩阵的那些分量。于是我们有
\begin{equation*}
\frac{d^2x^i}{d\tau^2} = \frac{1}{2}\left(\frac{dt}{d\tau}\right)^2\partial_i h_{00}.
\tag{4.17}
\end{equation*}
两边除以 $(dt/d\tau)^2$，效果是把左端的导数从 $\tau$ 换成 $t$，于是得到
\begin{equation*}
\frac{d^2x^i}{dt^2} = \frac{1}{2}\partial_i h_{00}.
\tag{4.18}
\end{equation*}
这看上去已经与 Newton 引力理论非常相像了。事实上，只要把此方程与 (4.1) 比较，就会发现二者相同，只需认定
\begin{equation*}
h_{00} = -2\Phi,
\tag{4.19}
\end{equation*}
换言之也就是
\begin{equation*}
g_{00} = -(1+2\Phi).
\tag{4.20}
\end{equation*}
因此我们已经证明：只要度规取 (4.20) 的形式，时空曲率确实足以在 Newton 极限下描述引力。当然，尚待完成的是找出度规的场方程——由它可推出度规确实取这种形式，而且对单个引力源天体能重新得到 Newton 公式
\begin{equation*}
\Phi = -\frac{GM}{r},
\tag{4.21}
\end{equation*}
不过这一点很快就会讲到。

我们上面概述的把物理定律推广到弯曲时空的直接流程确实有一些微妙之处，我们将在 4.7 节讨论。但就当前目的而言它已经绰绰有余，所以我们不要拖延，马上去完成任务的后一半：得到广义相对论中度规的场方程。

\section{Einstein 方程}

正如麦克斯韦方程组支配电场和磁场如何对电荷与电流作出响应，Einstein 场方程支配度规如何对能量和动量作出响应。归根结底，场方程必须靠假定提出、并接受实验检验，而不能从任何基石性原理导出；不过，我们可以基于合理性论证来为它提供动机。我们实际将采用两种方式：先通过一些非正式的类比推理——接近 Einstein 本人当年的思路——然后再从作用量出发，导出相应的运动方程。

非正式的论证从这样一个认识开始：我们想找一个方程来取代 Newton 势的 Poisson 方程：
\begin{equation*}
\nabla^2\Phi = 4\pi G\rho,
\tag{4.22}
\end{equation*}
其中 $\nabla^2=\delta^{ij}\partial_i\partial_j$ 是空间中的拉普拉斯算子（Laplacian），$\rho$ 是质量密度。［(4.21) 给出的 $\Phi$ 的显式形式是 (4.22) 在点状质量分布情形下的一个解。］我们苦苦寻觅的方程应当具备哪些特征？在 (4.22) 的左端，是一个作用在引力势上的二阶微分算子；右端则是对质量分布的一种量度。相对论性推广理应取张量之间方程的形式。我们知道质量密度的张量推广是什么：就是能动张量（energy-momentum tensor）$T_{\mu\nu}$。而引力势则应当被度规张量取代，因为在 (4.20) 中，我们正是靠把度规的微扰与 Newton 势联系起来才成功重现了引力。因此我们或许会猜测，新方程将把 $T_{\mu\nu}$ 设为正比于某个对度规导数而言为二阶的张量；大致像
\begin{equation*}
[\nabla^2 g]_{\mu\nu} \propto T_{\mu\nu},
\tag{4.23}
\end{equation*}
这个样子，当然我们要的是完全张量化的形式。

(4.23) 的左端并不是一个有意义的张量；它只是一个示意性记号，表明我们想要一个对称的 $(0,2)$ 张量，且对度规的导数是二阶的。第一个候选或许是让 d'Alembert 算符（d'Alembertian）$\Box=\nabla^\mu\nabla_\mu$ 作用在度规 $g_{\mu\nu}$ 上，但由度规相容性（metric compatibility）它自动为零。幸运的是，有一个显而易见的量，它不为零，而且由度规的二阶导数（以及一阶导数）构成：Riemann 张量 $R^\rho{}_{\sigma\mu\nu}$。回忆一下，Riemann 张量由 Christoffel 符号及其一阶导数构成，而 Christoffel 符号又由度规及其一阶导数构成，所以 $R^\rho{}_{\sigma\mu\nu}$ 含有 $g_{\mu\nu}$ 的二阶导数。它的指标数目不对，但我们可以把它缩并成指标数目恰当的 Ricci 张量 $R_{\mu\nu}$（而且还附赠对称性）。于是我们很自然会猜，引力场方程是
\begin{equation*}
R_{\mu\nu} = \kappa T_{\mu\nu},
\tag{4.24}
\end{equation*}
其中 $\kappa$ 为某个常数。事实上，Einstein 确实一度提出过这个方程。不幸的是，它存在一个与能量守恒有关的问题。如果我们想保持
\begin{equation*}
\nabla^\mu T_{\mu\nu} = 0,
\tag{4.25}
\end{equation*}
那么由 (4.24) 就会得到
\begin{equation*}
\nabla^\mu R_{\mu\nu} = 0.
\tag{4.26}
\end{equation*}
这在任意几何中显然不会成立；由 Bianchi 恒等式（3.150）我们已经看到
\begin{equation*}
\nabla^\mu R_{\mu\nu} = \frac{1}{2}\nabla_\nu R.
\tag{4.27}
\end{equation*}
但我们提议的场方程蕴涵 $R=\kappa g^{\mu\nu}T_{\mu\nu}=\kappa T$，把两者放在一起就得到
\begin{equation*}
\nabla_\mu T = 0.
\tag{4.28}
\end{equation*}
标量的协变导数就是偏导数，所以 (4.28) 是在告诉我们：$T$ 在整个时空中为常数。这极不可信，因为真空中 $T=0$，而物质中 $T\neq 0$。我们得再加把劲。

当然我们不必太费劲，因为我们已经知道一个由 Ricci 张量构造出来的、自动守恒的对称 $(0,2)$ 张量：Einstein 张量
\begin{equation*}
G_{\mu\nu} = R_{\mu\nu} - \frac{1}{2}Rg_{\mu\nu},
\tag{4.29}
\end{equation*}
它总是满足 $\nabla^\mu G_{\mu\nu}=0$。于是我们被引向提议
\begin{equation*}
G_{\mu\nu} = \kappa T_{\mu\nu}
\tag{4.30}
\end{equation*}
作为度规的场方程。（实际上，把 $R_{\mu\nu}-\frac{1}{2}Rg_{\mu\nu}$ 完整写出，可能比使用缩写 $G_{\mu\nu}$ 更常见。）这个方程满足所有显而易见的要求：右端把能量与动量密度写成一个对称、守恒的 $(0,2)$ 张量的协变表达式，而左端则是由度规及其一阶、二阶导数构造的对称、守恒的 $(0,2)$ 张量。剩下的只是确定比例常数 $\kappa$，并看看结果是否真的重现我们所知道的引力。换句话说，这个方程在 Newton 极限下能否预言引力势的 Poisson 方程？

为回答这个问题，注意把 (4.30) 两边缩并（在四维时空中）给出
\begin{equation*}
R = -\kappa T,
\tag{4.31}
\end{equation*}
利用它可以把 (4.30) 改写为
\begin{equation*}
R_{\mu\nu} = \kappa\left(T_{\mu\nu} - \frac{1}{2}Tg_{\mu\nu}\right).
\tag{4.32}
\end{equation*}
这是同一个方程，只是写法略有不同。我们想看看，它在弱场、时间无关、粒子缓慢运动的极限下是否预言 Newton 引力。我们考虑一个理想流体（perfect fluid）形式的能量-动量源，它满足
\begin{equation*}
T_{\mu\nu} = (\rho+p)U_\mu U_\nu + pg_{\mu\nu},
\tag{4.33}
\end{equation*}
其中 $U^\mu$ 是流体的四速度，$\rho$ 与 $p$ 是静止系的能量密度和动量密度。实际上，在 Newton 极限下我们可以忽略压强；粗略地说，当组成物体的粒子以接近光速的速度运动时，压强才变得重要，而按假设 Newton 极限不包含这种情形。所以我们实际考虑的是尘埃（dust）的能动张量：
\begin{equation*}
T_{\mu\nu} = \rho U_\mu U_\nu.
\tag{4.34}
\end{equation*}
我们考虑的“流体”是某个有质量的天体，比如地球或太阳。我们将在流体静止系（rest frame）中工作，在该系中
\begin{equation*}
U^\mu = (U^0, 0, 0, 0).
\tag{4.35}
\end{equation*}
类时分量可以借助归一化条件 $g_{\mu\nu}U^\mu U^\nu=-1$ 来确定。在弱场极限下，按照 (4.12) 与 (4.13)，我们写出
\begin{align*}
g_{00} &= -1 + h_{00}, \\
g^{00} &= -1 - h_{00}.
\tag{4.36}
\end{align*}
于是准确到 $h_{\mu\nu}$ 的一阶，我们得到
\begin{equation*}
U^0 = 1 + \frac{1}{2}h_{00}.
\tag{4.37}
\end{equation*}
不过，这实际上是多余的谨慎，因为我们反正要把四速度代进 (4.34)，而能量密度 $\rho$ 本来就已被当成小量（当 $\rho$ 趋于零时时空将是平直的）。所以在我们的近似水平上，可以直接取 $U^0=1$，相应地 $U_0=-1$。于是
\begin{equation*}
T_{00} = \rho,
\tag{4.38}
\end{equation*}
其他所有分量都为零。在这个极限下，静止能量 $\rho=T_{00}$ 将远大于 $T_{\mu\nu}$ 中的其他项，所以我们要聚焦于 (4.32) 的 $\mu=0$、$\nu=0$ 分量。迹在最低非平凡阶为
\begin{equation*}
T = g^{00}T_{00} = -T_{00} = -\rho.
\tag{4.39}
\end{equation*}

把它代人所提议的引力场方程 (4.32) 的 00 分量，得到
\begin{equation*}
R_{00} = \frac{1}{2}\kappa\rho.
\tag{4.40}
\end{equation*}
这是一个把度规的导数与能量密度联系起来的方程。要得到用度规表出的显式表达式，我们需要计算 $R_{00}=R^\lambda{}_{0\lambda 0}$。实际上只需要 $R^i{}_{0i0}$，因为 $R^0{}_{000}=0$。我们有
\begin{equation*}
R^i{}_{0j0} = \partial_j\Gamma^i_{00} - \partial_0\Gamma^i_{j0} + \Gamma^i_{j\lambda}\Gamma^\lambda_{00} - \Gamma^i_{0\lambda}\Gamma^\lambda_{j0}.
\tag{4.41}
\end{equation*}
这里第二项是时间导数，对静态场为零。第三、四项形如 $(\Gamma)^2$，由于 $\Gamma$ 是度规微扰的一阶小量，它们只在二阶才有贡献，可以忽略。于是剩下 $R^i{}_{0j0}=\partial_j\Gamma^i_{00}$。由此得到
\begin{align*}
R_{00} &= R^i{}_{0i0} \\
&= \partial_i\left[\frac{1}{2}g^{i\lambda}\left(\partial_0 g_{\lambda 0} + \partial_0 g_{0\lambda} - \partial_\lambda g_{00}\right)\right] \\
&= -\frac{1}{2}\delta^{ij}\partial_i\partial_j h_{00} \\
&= -\frac{1}{2}\nabla^2 h_{00}.
\tag{4.42}
\end{align*}
与 (4.40) 比较，我们看到 (4.30) 的 00 分量在 Newton 极限下预言
\begin{equation*}
\nabla^2 h_{00} = -\kappa\rho.
\tag{4.43}
\end{equation*}
由于 (4.19) 设定了 $h_{00}=-2\Phi$，只要取 $\kappa=8\pi G$，这正是 Poisson 方程 (4.22)。

看来我们的猜测 (4.30) 成功了。选定了能正确恢复 Newton 极限的归一化之后，我们就可以给出广义相对论的\textbf{Einstein 方程}：
\begin{equation*}
\boxed{\,R_{\mu\nu} - \frac{1}{2}Rg_{\mu\nu} = 8\pi G T_{\mu\nu}.\,}
\tag{4.44}
\end{equation*}
它告诉我们，时空曲率如何对能量-动量的存在作出反应。$G$ 当然是 Newton 引力常数，它与 $G_{\mu\nu}$ 的迹毫无关系。你可能听说过，Einstein 觉得方程左端漂亮而几何化，右端则多少不那么令人信服。

把 Einstein 方程改写成略有不同的形式有时很有用。按照 (4.31) 与 (4.32) 的做法，取 (4.44) 的迹可得 $R=-8\pi GT$。代入并把迹项移到右端，我们得到
\begin{equation*}
\boxed{\,R_{\mu\nu} = 8\pi G\left(T_{\mu\nu} - \frac{1}{2}Tg_{\mu\nu}\right).\,}
\tag{4.45}
\end{equation*}
它与 (4.44) 的差别纯属表面文章；实质上二者完全相同。我们常常对真空中的 Einstein 方程感兴趣，此时 $T_{\mu\nu}=0$（例如在恒星或行星之外）。当然，这时 (4.45) 的右端为零。因此真空 Einstein 方程就是
\begin{equation*}
\boxed{\,R_{\mu\nu} = 0.\,}
\tag{4.46}
\end{equation*}
它看上去稍不那么吓人，而且在物理上也相当有用。

\section{拉格朗日表述}

通往 Einstein 方程的另一条途径是最小作用量原理（principle of least action），正如我们在第 1 章末尾对平直时空中的经典场论所讨论的那样。让我们花点时间把那些结果推广到弯曲时空，然后看看什么样的拉格朗日量适合广义相对论。我们将在 $n$ 维中工作，因为我们的结果不依赖于维度；不过我们要假定度规具有洛伦兹号差。

考虑一个场论，其动力学变量是一组场 $\Phi^i(x)$。这样一个理论的经典解将是某个作用量 $S$ 的临界点，作用量一般表示为拉格朗日密度（Lagrange density）$\mathcal{L}$ 对空间的积分：
\begin{equation*}
S = \int \mathcal{L}(\Phi^i, \nabla_\mu\Phi^i)\, d^nx.
\tag{4.47}
\end{equation*}
注意，我们现在设想拉格朗日量是场及其协变导数（而非偏导数）的函数，这在弯曲空间中才是恰当的。还要注意，由于 $d^nx$ 是一个密度而不是张量，$\mathcal{L}$ 也是一个密度（因为二者的乘积必须是良定义的张量）；我们通常写成
\begin{equation*}
\mathcal{L} = \sqrt{-g}\,\widehat{\mathcal{L}},
\tag{4.48}
\end{equation*}
其中 $\widehat{\mathcal{L}}$ 确实是一个标量。你可能觉得明智的做法是忘掉我们称之为 $\mathcal{L}$ 的东西，只关注 $\widehat{\mathcal{L}}$；但事实上这两个量在不同场合各有用处。凡是对度规本身作变分时，起作用的都是 $\mathcal{L}$。与之相联系的 Euler--Lagrange 方程用到标量 $\widehat{\mathcal{L}}$，除把偏导数换成协变导数之外，与平直空间中的方程别无二致：
\begin{equation*}
\frac{\partial\widehat{\mathcal{L}}}{\partial\Phi} - \nabla_\mu\left(\frac{\partial\widehat{\mathcal{L}}}{\partial(\nabla_\mu\Phi)}\right) = 0.
\tag{4.49}
\end{equation*}
在推导这些方程时，我们用到 Stokes 定理（3.35）：
\begin{equation*}
\int_\Sigma \nabla_\mu V^\mu\sqrt{|g|}\, d^nx = \int_{\partial\Sigma} n_\mu V^\mu\sqrt{|\gamma|}\, d^{n-1}x,
\tag{4.50}
\end{equation*}
并把变分在无限远（边界处）取为零。于是分部积分的形式为
\begin{equation*}
\int A^\mu(\nabla_\mu B)\sqrt{-g}\, d^nx = -\int (\nabla_\mu A^\mu)B\sqrt{-g}\, d^nx + \text{边界项}.
\tag{4.51}
\end{equation*}
例如，第 1 章中考虑过的单标量场 $\phi$ 的作用量，其弯曲时空推广为
\begin{equation*}
S_\phi = \int\left[-\frac{1}{2}g^{\mu\nu}(\nabla_\mu\phi)(\nabla_\nu\phi) - V(\phi)\right]\sqrt{-g}\, d^nx,
\tag{4.52}
\end{equation*}
由此导出的运动方程为
\begin{equation*}
\Box\phi - \frac{dV}{d\phi} = 0,
\tag{4.53}
\end{equation*}
其中协变 d'Alembert 算符是
\begin{equation*}
\Box = \nabla^\mu\nabla_\mu = g^{\mu\nu}\nabla_\mu\nabla_\nu.
\tag{4.54}
\end{equation*}
正如在平直时空中一样，组合 $g^{\mu\nu}(\nabla_\mu\phi)(\nabla_\nu\phi)$ 常被缩写为 $(\nabla\phi)^2$。当然，协变导数作用在标量上时等价于偏导数，但明智的做法仍然是采用 $\nabla_\mu$ 记号；你说不定什么时候就要做分部积分，突然作用到一个矢量上。

热身之后，我们来为广义相对论构造作用量。现在的动力学变量是度规 $g_{\mu\nu}$；我们能用度规造出哪些标量来充当拉格朗日量呢？由于我们知道在任意一点可以把度规化为规范形式、并把度规的一阶导数化为零，任何非平凡的标量必然至少涉及度规的二阶导数。Riemann 张量当然是由度规的二阶导数构成的，而且我们此前论证过，由 Riemann 张量能构造出的唯一独立标量就是 Ricci 标量 $R$。我们虽然未曾证明、但事实如此的是：任何由度规与其一阶、二阶导数的乘积构成的非平凡张量，都可以用度规和 Riemann 张量表示出来。因此，由度规构造的、导数阶数不超过二阶的独立标量\emph{只有} Ricci 标量这一个。Hilbert 想必认为这就是拉格朗日量最简单的可能选择，于是提出
\begin{equation*}
S_H = \int\sqrt{-g}\, R\, d^nx,
\tag{4.55}
\end{equation*}
即\textbf{Hilbert 作用量}（Hilbert action，有时也称 Einstein--Hilbert 作用量）。正如我们将看到的，他是对的。

运动方程应当来自对作用量关于度规的变分。遗憾的是，这个作用量并不完全具有 (4.47) 的形式，因为它无法用 $g_{\mu\nu}$ 的协变导数写出来（协变导数会干脆为零）。因此，我们不再直接代入 Euler--Lagrange 方程，而是直接考察 $S_H$ 在度规小变分下的行为。实际上，关于逆度规（inverse metric）$g^{\mu\nu}$ 作变分更方便。由于 $g^{\mu\lambda}g_{\lambda\nu}=\delta^\mu_\nu$，而 Kronecker $\delta$ 在任何变分下都不变，度规变分与逆度规变分很容易互相表出：
\begin{equation*}
\delta g_{\mu\nu} = -g_{\mu\rho}g_{\nu\sigma}\,\delta g^{\rho\sigma},
\tag{4.56}
\end{equation*}
于是关于 $g^{\mu\nu}$ 变分的驻点等价于关于 $g_{\mu\nu}$ 变分的驻点。利用 $R=g^{\mu\nu}R_{\mu\nu}$，我们有
\begin{equation*}
\delta S_H = (\delta S)_1 + (\delta S)_2 + (\delta S)_3,
\tag{4.57}
\end{equation*}
其中
\begin{align*}
(\delta S)_1 &= \int d^nx\,\sqrt{-g}\, g^{\mu\nu}\delta R_{\mu\nu} \\
(\delta S)_2 &= \int d^nx\,\sqrt{-g}\, R_{\mu\nu}\delta g^{\mu\nu} \\
(\delta S)_3 &= \int d^nx\, R\,\delta\sqrt{-g}.
\tag{4.58}
\end{align*}
第二项 $(\delta S)_2$ 已经是某个表达式乘以 $\delta g^{\mu\nu}$ 的形式；让我们更仔细地考察其余两项。

回忆 Ricci 张量是 Riemann 张量的缩并，Riemann 张量由
\begin{equation*}
R^\rho{}_{\mu\lambda\nu} = \partial_\lambda\Gamma^\rho_{\nu\mu} + \Gamma^\rho_{\lambda\sigma}\Gamma^\sigma_{\nu\mu} - (\lambda\leftrightarrow\nu)
\tag{4.59}
\end{equation*}
给出。Riemann 张量关于度规的变分可以这样求：先求联络关于度规的变分，再代入这个表达式。不过，让我们考虑联络的任意变分，即作替换
\begin{equation*}
\Gamma^\rho_{\nu\mu} \rightarrow \Gamma^\rho_{\nu\mu} + \delta\Gamma^\rho_{\nu\mu}.
\tag{4.60}
\end{equation*}
变分 $\delta\Gamma^\rho_{\nu\mu}$ 是两个联络之差，因而本身是一个张量。于是我们可以取它的协变导数，
\begin{equation*}
\nabla_\lambda(\delta\Gamma^\rho_{\nu\mu}) = \partial_\lambda(\delta\Gamma^\rho_{\nu\mu}) + \Gamma^\rho_{\lambda\sigma}\delta\Gamma^\sigma_{\nu\mu} - \Gamma^\sigma_{\lambda\nu}\delta\Gamma^\rho_{\sigma\mu} - \Gamma^\sigma_{\lambda\mu}\delta\Gamma^\rho_{\nu\sigma}.
\tag{4.61}
\end{equation*}
此处以及后文，协变导数都是关于 $g_{\mu\nu}$ 取的，而不是关于 $g_{\mu\nu}+\delta g_{\mu\nu}$。有了这个表达式，再加上一点小小的劳动，容易证明：准确到变分的一阶，
\begin{equation*}
\delta R^\rho{}_{\mu\lambda\nu} = \nabla_\lambda(\delta\Gamma^\rho_{\nu\mu}) - \nabla_\nu(\delta\Gamma^\rho_{\lambda\mu}).
\tag{4.62}
\end{equation*}
欢迎你自己核对一下。于是，(4.58) 中第一项对 $\delta S$ 的贡献可以写成
\begin{align*}
(\delta S)_1 &= \int d^nx\sqrt{-g}\; g^{\mu\nu}\left[\nabla_\lambda(\delta\Gamma^\lambda_{\nu\mu}) - \nabla_\nu(\delta\Gamma^\lambda_{\lambda\mu})\right] \\
&= \int d^nx\sqrt{-g}\; \nabla_\sigma\left[g^{\mu\nu}(\delta\Gamma^\sigma_{\mu\nu}) - g^{\mu\sigma}(\delta\Gamma^\lambda_{\lambda\mu})\right],
\tag{4.63}
\end{align*}
其中我们用到了度规相容性，并重新命名了若干哑指标。现在可以代入 $\delta\Gamma^\sigma_{\mu\nu}$ 用 $\delta g^{\mu\nu}$ 表出的表达式，算出来是
\begin{equation*}
\delta\Gamma^\sigma_{\mu\nu} = -\frac{1}{2}\left[g_{\lambda\mu}\nabla_\nu(\delta g^{\lambda\sigma}) + g_{\lambda\nu}\nabla_\mu(\delta g^{\lambda\sigma}) - g_{\mu\alpha}g_{\nu\beta}\nabla^\sigma(\delta g^{\alpha\beta})\right],
\tag{4.64}
\end{equation*}
由此得到
\begin{equation*}
(\delta S)_1 = \int d^nx\sqrt{-g}\; \nabla_\sigma\left[g_{\mu\nu}\nabla^\sigma(\delta g^{\mu\nu}) - \nabla_\lambda(\delta g^{\sigma\lambda})\right],
\tag{4.65}
\end{equation*}
这也欢迎你自己核对。但 (4.63)［或 (4.65)］是一个矢量的协变散度对自然体积元的积分；由 Stokes 定理，它等于无限远处的边界贡献，只要让变分在无限远处为零，就可以把它取为零。因此这一项对总变分没有任何贡献。不过老实说，我们这里作了弊。边界项不仅包含度规的变分，还包含它的一阶导数，而后者按惯例并不取为零。就当前目的而言这无关紧要，但原则上我们可能在意边界上发生的事，那就必须在作用量中额外加进一项来处理这个微妙之处。

为了处理 $(\delta S)_3$ 项，我们需要用到下面这个事实，它对任何行列式非零的方阵 $M$ 都成立：
\begin{equation*}
\ln(\det M) = \mathrm{Tr}(\ln M).
\tag{4.66}
\end{equation*}
这里 $\ln M$ 由 $\exp(\ln M)=M$ 定义。对普通的数来说这显而易见，对矩阵则稍不那么直截了当。对这个恒等式取变分给出
\begin{equation*}
\frac{1}{\det M}\,\delta(\det M) = \mathrm{Tr}(M^{-1}\delta M).
\tag{4.67}
\end{equation*}
